% ==========================================
% IV. Experiments and Discussion (~1.65 pages)
% Logic: Datasets -> Baselines -> Metrics -> Results -> Discussion
% Key question: How do you demonstrate its effectiveness?
% ==========================================
\section{Experiments and Discussion}
\label{sec:experiments}

\subsection{Datasets}
Describe the data source, collection period, number of samples, features,
and train/test split.

\subsection{Baselines}
List the methods compared against the proposed scheme, e.g., linear
regression and random forest~\cite{breiman2001rf}.

\subsection{Evaluation Metrics}
The prediction performance is measured by the root mean square error (RMSE):
\begin{equation}
    \mathrm{RMSE} = \sqrt{\frac{1}{N} \sum_{i=1}^{N} \left( y_i - \hat{y}_i \right)^{2}},
    \label{eq:rmse}
\end{equation}
where $y_i$ and $\hat{y}_i$ are the actual and predicted values, respectively.

\subsection{Results}
Table~\ref{tab:comparison} compares the proposed scheme with the baselines.
Fig.~\ref{fig:results} visualizes the prediction results.

% ==========================================
% Comparison table (included by sections/04_experiments.tex)
% IEEE style: table captions go ABOVE the table.
% ==========================================
\begin{table}[htbp]
    \centering
    \caption{Performance Comparison of Different Methods}
    \label{tab:comparison}
    \begin{tabular}{lccc}
        \toprule
        Method & RMSE $\downarrow$ & MAE $\downarrow$ & $R^{2}$ $\uparrow$ \\
        \midrule
        Linear Regression & 0.000 & 0.000 & 0.000 \\
        Random Forest~\cite{breiman2001rf} & 0.000 & 0.000 & 0.000 \\
        Proposed Scheme & \textbf{0.000} & \textbf{0.000} & \textbf{0.000} \\
        \bottomrule
    \end{tabular}
\end{table}


\begin{figure*}[htbp]
    \centering
    \includegraphics[width=0.48\textwidth]{example-image-b}
    \hfill
    \includegraphics[width=0.48\textwidth]{example-image-c}
    \caption{Experimental results: (left) actual versus predicted values;
    (right) distribution of prediction errors.}
    \label{fig:results}
\end{figure*}

\subsection{Discussion}
Explain why the proposed scheme performs better (or worse) in certain cases,
analyze the limitations, and discuss practical implications.
